\VerdierPhysicalPageAnchor{0075}%
\VerdierSemanticAnchor{ch1-2-4-2-continuation}
Ces conventions s'appliquent en particulier à la catégorie
$\prod_{i\in[n]}\operatorname{comp}(C_i)$.

\VerdierSemanticAnchor{ch1-2-4-2-graded-product-category}
Remarquons que la catégorie
$\prod_{i\in[n]}\operatorname{comp}(C_i)$ est munie d'une structure de
$\mathbb Z[n]$-catégorie (1.3.11). Pour tout
$\sigma=\sum_{i\in[n]}l_i e_i$, $\sigma\in\mathbb Z[n]$, on a :

\VerdierSemanticAnchor{ch1-2-4-2-eq1}
\begin{equation}
  T(\sigma)(X_1,X_2,\ldots,X_n)
  = (X_1[l_1],X_2[l_2],\ldots,X_n[l_n]).
  \tag{2.4.2.1}
\end{equation}

\VerdierSemanticAnchor{ch1-2-4-2-summable-object}
Un objet $(X_1,\ldots,X_n)$ de
$\prod_{i\in[n]}\operatorname{comp}(C_i)$ est dit \emph{sommable} si,
pour tout $l\in\mathbb Z$, l'ensemble des
$\sigma=\sum_{i\in[n]}l_i e_i\in\mathbb Z[n]$ tels que
$\sum_{i\in[n]}l_i=l$ et $X_i^{l_i}\neq0$, pour au moins un
$i\in[n]$, est fini.

\VerdierSemanticAnchor{ch1-2-4-2-multiadditive-functor}
Soit maintenant
$F\colon\prod_{i\in[n]}C_i\longrightarrow C$ un foncteur multiadditif
(covariant). Pour tout objet $(X_1,\ldots,X_n)$ de
$\prod_{i\in[n]}\operatorname{comp}(C_i)$, posons :

\VerdierSemanticAnchor{ch1-2-4-2-eq2}
\begin{equation}
  F^{[n]}(X_1,\ldots,X_n)^\sigma
  =F(X_1^{l_1},\ldots,X_n^{l_n}),
  \tag{2.4.2.2}
\end{equation}
pour tout $\sigma=\sum_{i\in[n]}l_i e_i\in\mathbb Z[n]$,
\VerdierSemanticAnchor{ch1-2-4-2-eq2-differential}
\begin{equation*}
  d^{\sigma,i}_{F^{[n]}(X_1,\ldots,X_n)}
  =F(X_1^{l_1},\ldots,d_{X_i}^{l_i},\ldots,X_n^{l_n}),
  \qquad \forall i\in[n].
\end{equation*}

\VerdierSemanticAnchor{ch1-2-4-2-nfold-complex}
On a ainsi défini un complexe $n$-uple de $C$.

\VerdierSemanticAnchor{ch1-2-4-2-eq3}
Pour tout morphisme $(f_1,\ldots,f_n)$ de
$\prod_{i\in[n]}\operatorname{comp}(C_i)$, posons :
\begin{equation}
  F^{[n]}(f_1,\ldots,f_n)^\sigma
  =F(f_1^{l_1},\ldots,f_n^{l_n}),
  \tag{2.4.2.3}
\end{equation}
pour tout $\sigma=\sum_{i\in[n]}l_i e_i\in\mathbb Z[n]$.

\VerdierSemanticAnchor{ch1-2-4-2-morphism-of-complexes}
On a ainsi défini un morphisme de complexes $n$-uples.

\VerdierSemanticAnchor{ch1-2-4-2-eq4}
Ce qui précède définit visiblement un foncteur :
\begin{equation}
  F^{[n]}\colon
  \prod_{i\in[n]}\operatorname{comp}(C_i)
  \longrightarrow \operatorname{comp}^{n}(C).
  \tag{2.4.2.4}
\end{equation}

\VerdierSemanticAnchor{ch1-2-4-2-eq5}
Remarquons que, pour tout $\sigma\in\mathbb Z[n]$, on a :
\begin{equation}
  F^{[n]}T(\sigma)=T(\sigma)F^{[n]}.
  \tag{2.4.2.5}
\end{equation}

\VerdierSemanticAnchor{ch1-2-4-2-strict-graded-functor}
Par suite, le foncteur $F^{[n]}$ est un $\mathbb Z[n]$-foncteur strict.

\VerdierSemanticAnchor{ch1-2-4-2-preserves-summability}
Le foncteur $F^{[n]}$ transforme les objets sommables de
$\prod_{i\in[n]}\operatorname{comp}(C_i)$ en objets sommables de
$\operatorname{comp}^{n}(C)$.
